\documentclass[10pt,a4paper]{amsart}
\usepackage[margin=19mm]{geometry}
\usepackage{amsmath,amssymb,mathtools}
\usepackage[scr=rsfs,cal=euler]{mathalfa}
\usepackage{xfrac}
\usepackage[unicode,hidelinks,bookmarks=false]{hyperref}
\newcommand{\ie}{i.e.\ }
\newcommand{\cf}{cf.\ }
\newcommand{\resp}{respectively\ }
\newcommand{\Subsection}{\subsection}
\providecommand{\nobreakdash}{\nobreak\hbox{-}}
\setlength{\emergencystretch}{2em}
\multlinegap0pt
\usepackage{bm}
\usepackage{bbm}
\usepackage{dsfont}
\usepackage{xspace}
\usepackage{cancel}
\usepackage{mathtools}
\usepackage{amsthm}
\makeatletter
\@ifundefined{theorem}{\newtheorem{theorem}{Theorem}}{}
\@ifundefined{proposition}{\newtheorem{proposition}[theorem]{Proposition}}{}
\makeatother
\newcommand\restr[2]{{% we make the whole thing an ordinary symbol
  \left.\kern-\nulldelimiterspace % automatically resize the bar with \right
  #1 % the function
  %\vphantom{\big|} % pretend it's a little taller at normal size
  \right|_{#2} % this is the delimiter
}}
\newcommand{\Cinfty}{\mathscr{C}^\infty}
\newcommand{\R}{\mathbb{R}}
\newcommand{\T}{\mathrm{T}}
\newcommand{\bh}{{\boldsymbol{h}}}
\newcommand{\bomega}{{\boldsymbol{\omega}}}
\newcommand{\inn}{i}
\renewcommand{\d}{\mathrm{d}}
\title[An energy-momentum method for ordinary differential equation]{An energy-momentum method for ordinary differential equations with an underlying $k$-polysymplectic manifold}
\author{Leonardo Colombo, Javier de Lucas, Xavier Rivas, Bartosz M. Zawora}
\date{Source: arXiv:2311.15035v1 (CC BY 4.0)}
\begin{document}
\maketitle

\noindent\emph{Source and licence.} An energy-momentum method for ordinary differential equations with an underlying $k$-polysymplectic manifold. Version arXiv:2311.15035v1 (CC BY 4.0), distributed under the Creative Commons Attribution 4.0 International licence, as recorded in the arXiv licence metadata for this exact version and shown on the abstract page https://arxiv.org/abs/2311.15035v1. Full rights notice: licenses/arxiv-2311.15035-CC-BY-4.0.txt.\par\smallskip

\noindent\textit{Editorial scope.} Counted unit: the Proposition whose source label is \texttt{Prop::GaugeDir} in the article (source0.tex lines 1220--1228, source label \texttt{Prop::GaugeDir}), delivered together with the article's own complete original proof of it (source0.tex lines 1229--1253, one \texttt{proof} environment of 25 lines). No mathematics is written, repaired or re-derived: statement and proof are the article's own text line for line. Required context retained uncounted: Eq::SecVarf (lines 1195--1198). Every definition, display and label of those lines is retained unchanged. Omitted from this package and not redistributed: 1--1194, 1199--1219, 1254--1940. Nothing inside a retained excerpt is omitted or altered: every retained body is delivered exactly as the article's lines are written, so no omission receipt of any kind accompanies this package. Citations: the retained excerpts carry no citation command at all, so no bibliography entry of the article is transcribed and no reference is redistributed. Reference adaptation: none was needed and none was made; every reference inside the delivered unit resolves inside it, so no unresolved reference or citation is produced. Printed slips kept literal: no printed word, symbol, hypothesis or number of the counted statement or of its proof was altered. Layout and class: the article's own class is \texttt{article} with packages including inputenc, fontenc, hyperref, colortbl, graphicx, srcltx, amsmath, amssymb, bm, enumerate, amsthm, xy, mathrsfs, mathtools; the wrapper is the lane's pinned \texttt{amsart} editorial wrapper at 10pt on A4, which restates the theorem-like environments the retained text needs and adds the article's own macro definitions that the retained text needs (\texttt{\textbackslash Cinfty}, \texttt{\textbackslash R}, \texttt{\textbackslash T}, \texttt{\textbackslash bh}, \texttt{\textbackslash bomega}, \texttt{\textbackslash d}, \texttt{\textbackslash inn}, \texttt{\textbackslash restr}), with extra wrapper packages (bm, bbm, dsfont, xspace, cancel, mathtools). The delivered numbering is the unit's own. Assets: no retained span contains a figure environment, a graphics-inclusion command, a PDF-page inclusion, a table or a TikZ picture; no image file is copied into this package, the package declares no asset dependency and no image file is redistributed with it. Licence: version arXiv:2311.15035v1 of this article is distributed under the Creative Commons Attribution 4.0 International licence, as recorded in the arXiv licence metadata for this exact version and shown on the abstract page https://arxiv.org/abs/2311.15035v1. A full rights notice is delivered at licenses/arxiv-2311.15035-CC-BY-4.0.txt. Characters: every retained excerpt is pure ASCII, so the delivered engine, XeLaTeX with its default Latin Modern text font, reports no missing character.
\begin{equation}
\label{Eq::SecVarf}
\left(\delta^2 \boldsymbol{h}_\xi\right)_{z_e}(v_1,v_2)=\sum^k_{\alpha=1}\inn{Y}\left(\d\left(\inn{X}\d h^\alpha_\xi\right)\right)_{z_e}\otimes e_\alpha\,,
\end{equation}

\begin{proposition}
\label{Prop::GaugeDir}
     Let $(P,\bomega,\bh,\mathbf{J}^\Phi)$ be a $G$-invariant $\bomega$-Hamiltonian system and let $z_e\in P$ be a $k$-polysymplectic relative equilibrium point of $X_{\bh}$. Then,
    \[ (\delta^2\boldsymbol{h}_\xi)_{z_e}((\zeta_P)_{z_e},v_{z_e})=0\,,\qquad \forall \zeta\in \mathfrak{g}^{\boldsymbol{\Delta}}_{\boldsymbol{\mu}_e}\,,\qquad \forall v_{z_e}\in \T_{z_e}({\bf J}^{\Phi-1}({\bm \mu}_e))\,,
    \]
    with ${\bm \mu}_e={\bf J}^{\Phi}(z_e)$. Moreover,
     \begin{equation}\label{eq:SecDeg} (\delta^2{h}^\alpha_\xi)_{z_e}(Y_{z_e},\cdot)=0\,,\qquad \forall \,Y_{z_e}\in \ker \omega_{z_e}^\alpha\cap \T_{z_e}({\bf J}^{\Phi-1}(\bm\mu_e))\,,\qquad \alpha=1,\ldots,k\,.
    \end{equation}
\end{proposition}

\begin{proof}
First, since $\bh\in \Cinfty(P,\R^k)$ is $G$-invariant and $\mathbf{J}^\Phi$ is equivariant with respect to the $k$-polysymplectic affine Lie group action $\boldsymbol{\Delta}:G \times (\mathfrak{g}^*)^k\rightarrow (\mathfrak{g}^*)^k$, then for every $g\in G$ and $p\in P$, one has
\begin{multline*}
\boldsymbol{h}_\xi(\Phi_g(p)) = \boldsymbol{h}(\Phi_g(p))-\langle \mathbf{J}^\Phi(\Phi_g(p)),\xi\rangle +\langle {\bm \mu}_e,\xi\rangle \\
= \boldsymbol{h}(p)-\langle \boldsymbol{\Delta}_g\mathbf{J}^\Phi(p),\xi\rangle +\langle \bm\mu_e,\xi\rangle %&=\boldsymbol{h}(p)-\sum_{\alpha=1}^k\langle \mathbf{J}_\alpha^\Phi(p),\boldsymbol{\Delta}_{\alpha g}^T\xi\rangle \otimes e_\alpha+\langle \bm\mu_e,\xi\rangle \,,\\
=\boldsymbol{h}(p)-\sum_{\alpha=1}^k\langle \mathbf{J}_\alpha^\Phi(p),\Delta_{g\alpha}^T\xi\rangle \otimes e_\alpha+\langle \bm\mu_e,\xi\rangle,
\end{multline*}
where $\boldsymbol{\Delta}^T_g:\mathfrak{g}^k\rightarrow \mathfrak{g}^k$ is the transpose of $\boldsymbol{\Delta}_g$ for $g\in G$ and $\Delta_{g1},\ldots,\Delta_{gk}$  are its components. Let us substitute $g=\exp(t\zeta)$, with $\zeta\in\mathfrak{g}$, and differentiate with respect to $t$. Then,
\begin{equation}
\label{Eq::1stvar}
\left(\inn{\zeta_P}\d\boldsymbol{h}_\xi\right)_{z_e}=-\sum_{\alpha=1}^k\left\langle \mathbf{J}_\alpha^\Phi(p),\restr{\frac{\d}{\d t}}{t=0}{\Delta}^T_{\exp{(t\zeta)\alpha}}\xi \right\rangle\otimes e_\alpha=-\sum_{\alpha=1}^k\left\langle \mathbf{J}_\alpha^\Phi(p),(\zeta^{{\Delta}_\alpha}_{\mathfrak{g}})_{\xi}\right\rangle \otimes e_\alpha\,,
\end{equation}
where $(\zeta^{{\Delta}_\alpha}_{\mathfrak{g}})_{\xi}$ is the fundamental vector field of ${\Delta}_\alpha^T:G\times\mathfrak{g}\rightarrow \mathfrak{g}$ at $\xi\in\mathfrak{g}$ for $\alpha=1,\ldots,k$. Taking the second variation of \eqref{Eq::1stvar} relative to $p\in P$, evaluating at $z_e\in P$, and contracting with $v_{z_e}$, one has
\[
\left(\delta^2 \boldsymbol{h}_\xi \right)_{z_e}((\zeta_P)_{z_e},v_{z_e})=-\sum_{\alpha=1}^k\left\langle \T_{z_e}\mathbf{J}_\alpha^\Phi\left(v_{z_e}\right), (\zeta^{{\Delta}_\alpha}_{\mathfrak{g}})_{\xi}\right\rangle\otimes e_\alpha .
\]
%By Proposition \ref{Lemm::NonAdPerpPS}, 
Therefore, the second variation $\left(\delta^2 \boldsymbol{h}_\xi \right)_{z_e}((\zeta_P)_{z_e},v_{z_e})$ vanishes since $v_{z_e}\in \T_{z_e}(\mathbf{J}^{\Phi-1}(\boldsymbol{\mu}_e))\subset \ker \T_{z_e}\mathbf{J}^{\Phi}_\alpha$.

Concerning \eqref{eq:SecDeg}, it is a consequence of \eqref{Eq::SecVarf} and the fact that, for every vector field $Y$ on ${\bf J}^{\Phi-1}(\bm\mu_e)$ taking values in $\ker\omega^\alpha\cap \T({\bf J}^{\Phi-1}(\bm\mu_e))$, it follows that
$$
\iota_Y\d h^\alpha=\omega^\alpha(X_{\bm h},Y)=0\,,\qquad \iota_Y\d \langle {\bf J}_\alpha^\Phi,\xi\rangle =\omega^\alpha(\xi_P,Y)=0\,,
$$
for $\alpha=1,\ldots,k$ and every $\xi\in \mathfrak{g}$ on ${\bf J}^{\Phi-1}(\bm\mu_e)$. 
\end{proof}

\end{document}
